\newpage

\section{point 类}
在 tkz-elements 中，一个点在内部表示为一个复数（即它的附标）。这一选择使得平移、旋转、缩放等计算可以简洁地进行，同时在渲染时仍与 TikZ/tkz-euclide 兼容。


变量 \tkzVar{point}{z} 保存一个用于存储点的表。它是 \textbf{必需的}，由本宏包自动初始化（例如 \tkzname{ z = \{\}}）。


\vspace{1em}
\tkzClass{point} 类构成整个框架的基础。它是一个混合类，既表示平面中的点，也表示复数。其基本原理如下：

\begin{itemize}
  \item 平面配备了一个正交基（OCCS，见 [\ref{sec:orthonormal_cartesian_coordinate_system}]），使我们可以通过横坐标和纵坐标确定一个点的位置。
  \item 类似地，任何复数都可看作一个有序的实数对（其实部和虚部）。
  \item 因此，平面可以等同于复平面，复数 $x + iy$ 由平面中坐标为 $(x, y)$ 的点来表示。
\end{itemize}

因此，像 $A$ 这样的点被存储为对象 \code{z.A}，其坐标及相关属性都封装在该对象内部。

\vspace{1em}
\code{point} 对象具有以下几个属性：
\begin{mybox}{}
\begin{itemize}
  \item \code{re} --> 实部（横坐标），
  \item \code{im} --> 虚部（纵坐标），
  \item \code{type} --> 对象的类型（在此情形下始终为 \code{"point"}），
  \item \code{arg} --> 复数的辐角（相对于 x 轴的角度），
  \item \code{mod} --> 复数的模（到原点的距离）。
\end{itemize}
\end{mybox}



\subsection{创建点}
\label{sub:creating_a_point}

通过在当前正交笛卡尔坐标系（OCCS）中给出坐标来创建点。推荐的形式为：

\begin{mybox}
  \code{z.A = point(x, y)}
\end{mybox}

其中 \code{x} 和 \code{y} 是实数，对应于该点的 $x$ 和 $y$ 坐标。

\medskip

在内部，这会创建一个复数 $x + i y$，并以键 \code{"A"} 存储在表 \tkzVar{point}{z} 中。表 \tkzVar{point}{z} 用于按标签引用所有点。

\medskip

另外，也可以使用更显式的语法：

\begin{center}
  \code{z.A = point:new(x, y)}
\end{center}


两种形式等价。简写构造函数自第 4 版起提供，为便于阅读和一致起见推荐使用。


\def\size{42mm}
\begin{tikzpicture}[remember picture]
\node[ draw, fill=red!10] (tbl) {%
 \centering
\begin{minipage}{\size}

   \hspace{\fill} \texttt{属性}\hspace{\fill}

       \tikz\node[minimum width=\size,font=\small,
    draw, fill=cyan!10,
    rectangle split, rectangle split parts=6
  ] {
    \texttt{re (实部)}
    \nodepart{two}\texttt{im (虚部)}
    \nodepart{three}\texttt{type = 'point'}
    \nodepart{four}\texttt{argument (弧度)}
    \nodepart{five}\texttt{modulus (cm)}
    \nodepart{six}\texttt{mtx (矩阵)}
  };

    \hspace{\fill}  \texttt{方法}\hspace{\fill}

    \tikz\node[minimum width=\size,font=\small,
    draw, fill=orange!20,sharp corners,
    rectangle split, rectangle split parts=4
  ] {
   \texttt{homothety(coeff,obj)}
    \nodepart{two}\texttt{rotation(angle, object)}
    \nodepart{three}\texttt{symmetry (object)}
    \nodepart{four}\texttt{\ldots}
  };
\end{minipage}};
 \node[ draw, fill=red!10,,minimum height = 2em,
  rounded corners,anchor=south] (tc) at (tbl.north){类 |Point|};
\end{tikzpicture}
\hspace{5cm}\begin{tikzpicture}[remember picture]
   \node[ draw, fill=red!10] (tbl) {%
 \centering
\begin{minipage}{\size}
   \hspace{\fill}    \texttt{属性}\hspace{\fill}

        \tikz\node[minimum width=\size,font=\small,
    draw, fill=cyan!10,
    rectangle split, rectangle split parts=6
  ] {
    \texttt{re = 1}
    \nodepart{two}\texttt{im = 2}
    \nodepart{three}\texttt{type = 'point'}
    \nodepart{four}\texttt{argument = atan(2)}
     \nodepart{five}\texttt{modulus = $\sqrt{5}$}
      \nodepart{six}\texttt{mtx = \{\{1\},\{2\}\}}
  };

    \hspace{\fill}  \texttt{方法}\hspace{\fill}

        \tikz\node[minimum width=\size,font=\small,
    draw, fill=orange!20,sharp corners,
    rectangle split, rectangle split parts=4
  ] {
   \texttt{homothety(coeff,obj)}
    \nodepart{two}\texttt{rotation(angle, object)}
    \nodepart{three}\texttt{symmetry (object)}
    \nodepart{four}\texttt{\ldots}
  };
\end{minipage}
     };
 \node[ draw, fill=red!10,remember picture,minimum height = 2em,
  rounded corners,anchor=south] (to) at (tbl.north){对象 |z.A|};
\end{tikzpicture}

\begin{tikzpicture}[remember picture,overlay]
\draw [thick,->](tc.east) -- (to.west);
\end{tikzpicture}


\subsection{点的属性}
\label{sub:attributes_of_a_point}
% 方法 \tkzMeth{point}{new}

\begin{mybox}
   创建 \\
   |z.A = point(1, 2) |
\end{mybox}
 点 $A$ 的坐标为 $x=1$、$y=2$。如果你采用 |z.A| 记法，则 $A$ 将作为 \TIKZ\ 或 \tkzNamePack{tkz-euclide} 中的节点被引用。

这创建了一个坐标为 $1$、$2$ 且名为 $A$ 的固定点。记法 \tkzname{z.A} 表示坐标将存储在变量 \tkzVar{point}{z} 所对应的表中（呼应复数附标的记法），其中 $A$ 是点的名称，也是用于访问各值的键。

\label{point:attributes}
\bgroup
  \small
  \catcode`_=12
  \captionof{table}{点的属性。}
  \begin{tabular}{lll}
  \toprule
  \texttt{属性}      & \textbf{说明}&   \texttt{示例 / 参考} \\
  \midrule
  \tkzAttr{point}{type}     &  对象类型名，始终为 \texttt{"point"}   & \\
  \tkzAttr{point}{re}       &  实部（即 $x$ 坐标）      & [\ref{ssub:methods}] \\
  \tkzAttr{point}{im}       &  虚部（即 $y$ 坐标）    & [\ref{ssub:methods}]  \\
  \tkzAttr{point}{argument} &   附标的辐角（弧度角） & $\approx$ \code{0.785398...} \quad  [\ref{ssub:example_point_attributes}] \\
  \tkzAttr{point}{modulus}  & 附标的模（到原点的距离）  &  $\sqrt{5} \approx$ \code{2.2360...} \quad [\ref{ssub:example_point_attributes}] \\
  \tkzAttr{point}{mtx}  & 以列向量形式表示的矩阵   & \code{z.A.mtx = \{\{1\},\{2\}\}} \quad  [\ref{ssub:example_point_attributes}] \\
  \bottomrule
  \end{tabular}
\egroup


\subsubsection{示例：点的属性}
\label{ssub:example_point_attributes}

\directlua{
init_elements()
   z.M = point(1, 2)}
\hspace*{\fill}

\begin{verbatim}
\directlua{
   init_elements()
   z.M = point(1, 2)}
\end{verbatim}
\pgfkeys{/pgf/number format/.cd,std,precision=2}
\let\pmpn\pgfmathprintnumber
\DeleteShortVerb{\|}

\begin{verbatim}
\begin{tikzpicture}[scale = 1]
\pgfkeys{/pgf/number format/.cd,std,precision=2}
\let\pmpn\pgfmathprintnumber
\tkzDefPoints{2/4/M,2/0/A,0/0/O,0/4/B}
\tkzLabelPoints(O)
\tkzMarkAngle[fill=gray!30,size=1](A,O,M)
\tkzLabelAngle[pos=1,right](A,O,M){%
   $\theta \approx \pmpn{\tkzUseLua{z.M.argument}}$ rad}
\tkzDrawSegments(O,M)
\tkzLabelSegment[above,sloped](O,M){%
    $|z_M| =\sqrt{5}\approx \pmpn{%
   \tkzUseLua{z.M.modulus}}$ cm}
\tkzLabelPoint[right](M){$M: z_M = 1 + 2i$}
\tkzDrawPoints(M,A,O,B)
\tkzPointShowCoord(M)
\tkzLabelPoint[below,teal](A){$\tkzUseLua{z.M.re}$}
\tkzLabelPoint[left,teal](B){$\tkzUseLua{z.M.im}$}
\tkzDrawSegments[->,add = 0 and 0.25](O,B O,A)
\end{tikzpicture}
\end{verbatim}


\begin{center}
   \begin{tikzpicture}
   \pgfkeys{/pgf/number format/.cd,std,precision=2}
   \let\pmpn\pgfmathprintnumber
   \tkzDefPoints{2/4/M,2/0/A,0/0/O,0/4/B}
   \tkzLabelPoints(O)
   \tkzMarkAngle[fill=gray!30,size=1](A,O,M)
   \tkzLabelAngle[pos=1,right](A,O,M){%
   $\theta \approx \pmpn{\tkzUseLua{z.M.argument}}$ rad}
   \tkzDrawSegments(O,M)
   \tkzLabelSegment[above,sloped](O,M){%
   $|z_M| =\sqrt{5}\approx \pmpn{\tkzUseLua{z.M.modulus}}$ cm}
   \tkzLabelPoint[right](M){$M: z_M = 1 + 2i$}
   \tkzDrawPoints(M,A,O,B)
   \tkzPointShowCoord(M)
   \tkzLabelPoint[below,teal](A){$\tkzUseLua{z.M.re}$}
   \tkzLabelPoint[left,teal](B){$\tkzUseLua{z.M.im}$}
   \tkzDrawSegments[->,add = 0 and 0.25](O,B O,A)
   \begin{scope}[every annotation/.style={fill=lightgray!15,anchor = east}]
   \node [annotation,font =\small,text width=6cm] at (current bounding box.west) {
\texttt{z.M} 的属性
   \begin{mybox}{}
     \begin{itemize}
     \item \texttt{z.M.re} = 1
     \item \texttt{z.M.im} = 2
     \item \texttt{z.M.type} = point
     \item \texttt{z.M.argument} = $\theta \approx \pmpn{\tkzUseLua{z.M.argument}}$ rad
     \item \texttt{z.M.modulus} = $|z_M| =\sqrt{5}\approx \pmpn{\tkzUseLua{z.M.modulus}}$ cm
     \item \texttt{z.M.mtx} =  \tkzUseLua{z.M.mtx:print()}
     \end{itemize}
   \end{mybox}
      };
   \end{scope}
   \end{tikzpicture}
\end{center}

 \MakeShortVerb{\|}


\subsubsection{属性 \tkzAttr{point}{mtx}}
\label{ssub:attribute_iattr_point_mtx}

该属性使点能够与矩阵结合使用。

\vspace{1em}
\begin{minipage}{.5\textwidth}
\begin{verbatim}
  \directlua{
  z.A = point(2, -1)
  z.A.mtx:print()}
\end{verbatim}
\end{minipage}
\begin{minipage}{.5\textwidth}
\begin{center}
  \directlua{
  z.A = point(2, -1)
  z.A.mtx:print()}
\end{center}
\end{minipage}

\subsubsection{复平面（Argand 图）}
\label{ssub:argand_diagram}
\normalsize
\begin{minipage}{\textwidth}
\begin{verbatim}
\directlua{
  init_elements()
  z.A = point(2, 3)
  z.O = point(0, 0)
  z.I = point(1, 0)}
\begin{tikzpicture}
   \tkzGetNodes
   \tkzInit[xmin=-4,ymin=-4,xmax=4,ymax=4]
   \tkzDrawCircle[dashed,red](O,A)
   \tkzPointShowCoord(A)
   \tkzDrawPoint(A)
   \tkzLabelPoint[above right](A){\normalsize $a+ib$}
   \tkzDrawX\tkzDrawY
   \tkzDrawSegment(O,A)
   \tkzLabelSegment[above,anchor=south,sloped](O,A){ OA = modulus of $z_A$}
  \tkzLabelAngle[anchor=west,pos=.5](I,O,A){$\theta$ = argument of $z_A$}
\end{tikzpicture}
\end{verbatim}
\end{minipage}

\begin{minipage}{\textwidth}
\directlua{
 init_elements()
 z.A = point(2, 3)
 z.O = point(0, 0)
 z.I = point(1, 0)}
\begin{center}
\begin{tikzpicture}
  \tkzGetNodes
  \tkzInit[xmin=-4,ymin=-4,xmax=4,ymax=4]
  \tkzDrawCircle[dashed,red](O,A)
  \tkzPointShowCoord(A)
  \tkzDrawPoint(A)
  \tkzLabelPoint[above right](A){\normalsize $a+ib$}
  \tkzDrawX\tkzDrawY
  \tkzDrawSegment(O,A)
  \tkzLabelSegment[above,anchor=south,sloped](O,A){ OA = $z_A$ 的模}
  \tkzLabelAngle[anchor=west,pos=.5](I,O,A){$\theta$ = $z_A$ 的辐角}
\end{tikzpicture}
\end{center}
\end{minipage}

\newpage
\subsection{point 类的方法}

下表所列的方法都是标准方法，在本文档末尾的各示例中被广泛使用。这些方法每一个都返回一个 \tkzNameObj{point} 对象。

\medskip

关于使用复数和运算符重载的更高级运算，请参见第~\ref{sub:complex_numbers} 节，其中描述了可用的元方法。


\vspace{1em}

\bgroup
\catcode`_=12
\small
\captionof{table}{\code{class} point 的函数与方法。}\label{point:methods}
\begin{tabular}{lll}
\toprule

\texttt{方法} & \texttt{参考} \\
\midrule
\texttt{构造函数} &\\
\midrule
\tkzMeth{point}{new(r,r)} &[\ref{ssub:method_point_new}; \ref{ssub:method_normalize}]\\
\tkzMeth{point}{polar(d,an)} &[\ref{ssub:method_point_polar}] \\
\tkzMeth{point}{polar\_deg(d,an)} &[\ref{ssub:method_point_polar_deg}]   \\
\midrule
\texttt{返回实数的方法} & \\
\midrule
\tkzMeth{point}{get()}     &[\ref{ssub:method_point_get}] \\
\midrule
\texttt{返回点的方法} & \\
\midrule
\tkzMeth{point}{north(r)} & [\ref{ssub:methods}]   \\
\tkzMeth{point}{south(r)} &  \\
\tkzMeth{point}{east(r)}  &  \\
\tkzMeth{point}{west(r)}  &  \\
\tkzMeth{point}{normalize()}         & [\ref{ssub:method_normalize}] \\
\tkzMeth{point}{normalize\_from(pt)} & [\ref{ssub:method_normalize}] \\
\tkzMeth{point}{orthogonal(d)}       & [\ref{ssub:orthogonal_method}]\\
\tkzMeth{point}{at()}                &  [\ref{ssub:_point_at_method}] \\
\tkzMeth{point}{shift\_orthogonal\_to(pt, d)}& [\ref{shift_orthogonal}] \\
\tkzMeth{point}{shift\_collinear\_to(pt, d)} & [\ref{shift_collinear}] \\
\midrule
\texttt{返回圆的方法} & \\
\midrule
\tkzMeth{point}{PPP(a,b)} &[\ref{ssub:point_PPP}]\\
\midrule
\texttt{返回对象的方法} & \\
\midrule
\tkzMeth{point}{symmetry(obj)} & [\ref{ssub:object_symmetry}] \\
\tkzMeth{point}{rotation(an, obj)}  &  [\ref{ssub:object_rotation}] \\
\tkzMeth{point}{homothety(r,obj)}    & [\ref{ssub:method_point_homothety_k_obj}]   \\
\midrule

\texttt{实用工具} &\\
\midrule
\tkzMeth{point}{identity(pt)} & [\ref{ssub:identity}]\\
\tkzMeth{point}{print()} & [\ref{ssub:print}]\\
\bottomrule %
\end{tabular}
\egroup


\subsubsection{方法 \tkzMeth{point}{new(r, r)}}
\label{ssub:method_point_new}

该方法使用笛卡尔坐标在平面中创建一个点。\textbf{简写构造函数} \tkzMeth{point}{(r, r)} 自第 4 版起提供，为便于阅读和一致起见推荐使用。

\medskip

它接受两个实数作为参数：第一个表示 \emph{横坐标}（实部），第二个表示 \emph{纵坐标}（虚部）。在内部，该点被视为一个复数，并存储在全局表 \tkzname{z} 中。

\medskip

所得对象的类型为 \tkzNameObj{point}，可用于进一步的几何构造，或用 \tkzNamePack{tkz-euclide} 显示。

\medskip

\smallbf{注意：} 默认单位为厘米，与 \tkzNamePack{tkz-euclide} 的约定一致。所有坐标均在正交笛卡尔坐标系中解释。

\vspace{1em}

\begin{tkzexample}[latex = 7cm]
\directlua{
  init_elements()
  z.A  = point(0, 0)
  z.B  = point(2, 1)}
\begin{tikzpicture}[gridded]
  \tkzGetNodes
  \tkzDrawPoints(A,B)
  \tkzLabelPoints(A,B)
\end{tikzpicture}
\end{tkzexample}

\subsubsection{方法 \tkzMeth{point}{get()}}
\label{ssub:method_point_get}

\code{get} 方法用于取回一个点的笛卡尔坐标。

\medskip

设 $I$ 为两条直线的交点。你可以用两种等价方式获取它的坐标：
\begin{itemize}
  \item 直接使用该点的属性：
    \begin{center}
    \code{x = z.I.re} \quad 与 \quad \code{y = z.I.im}
    \end{center}
  \item 或使用 \code{get()} 方法：
    \begin{center}
    \code{x, y = z.I:get()}
    \end{center}
\end{itemize}


该方法提升了代码的可读性，也便于把坐标传给那些期望数值参数的函数。

\vspace{1em}

\begin{minipage}{.5\textwidth-1cm}
\directlua{
 init_elements()
 z.A, z.B = point(0, 0), point(5, -1)
 z.C, z.D = point(1, -4), point(4, 2)
 L.AB = line(z.A, z.B)
 L.CD = line(z.C, z.D)
 z.I = intersection(L.AB, L.CD)
 x, y = z.I:get()
 tex.print("$x_I = $"..x)
 tex.print('\\\\')
 tex.print("$y_I = $"..y)}
\end{minipage}
\begin{minipage}{.5\textwidth}
\begin{tkzexample}[code only]
\directlua{
 init_elements()
 z.A, z.B = point(0, 0), point(5, -1)
 z.C, z.D = point(1, -4), point(4, 2)
 L.AB = line(z.A, z.B)
 L.CD = line(z.C, z.D)
 z.I = intersection(L.AB, L.CD)
 x, y = z.I:get()
 tex.print("$x_I = $"..x)
 tex.print('\\\\')
 tex.print("$y_I = $"..y)}
\end{tkzexample}
\end{minipage}

\subsubsection{函数 \tkzFct{point}{polar(r, an)}}
\label{ssub:method_point_polar}

该方法使用极坐标在平面中创建一个点。

\medskip

它接受两个参数：
\begin{itemize}
  \item \code{r} —> 模（到原点的距离），
  \item \code{an} —> 辐角（弧度角）。
\end{itemize}


在内部，该点表示为复数：\code{r * exp(i * an)}。该方法特别适用于在圆上构造点，或用角度和距离定义点。

\medskip

\textbf{注意：} 默认单位为厘米，与 \tkzNamePack{tkz-euclide} 的约定一致。所有坐标均在正交笛卡尔坐标系中解释。

\vspace{1em}

\begin{mybox}
\begin{verbatim}
  z.B = polar(2, math.pi / 4)
  or
  z.B = point:polar(2, math.pi / 4)
\end{verbatim}
\end{mybox}

\begin{tkzexample}[latex = 7cm]
\directlua{
 init_elements()
 z.O = point(0, 0)
 z.A = point(3, 0)
 z.F = polar(3, math.pi / 3)}
\begin{center}
\begin{tikzpicture}[scale=.75]
   \tkzGetNodes
   \tkzDrawCircle(O,A)
   \tkzDrawSegments[new](O,A)
   \tkzDrawSegments[purple](O,F)
   \tkzDrawPoints(A,O,F)
   \tkzLabelPoints[below right=6pt](A,O)
   \tkzLabelPoints[above](F)
\end{tikzpicture}
\end{center}

\end{tkzexample}

\subsubsection{方法 \tkzMeth{point}{polar\_deg(d,an)}}
\label{ssub:method_point_polar_deg}

与上一个相同，只是角度以度数给出。

\begin{tkzexample}[latex=7cm]
\directlua{
 init_elements()
 z.A = point(0, 0)
 z.B = polar_deg(3, 60)
 z.C = polar_deg(3, 0)}
\begin{center}
  \begin{tikzpicture}[gridded]
  \tkzGetNodes
  \tkzDrawPolygon(A,B,C)
  \tkzDrawPoints(A,B,C)
  \tkzLabelPoints(A,C)
  \tkzLabelPoints[above](B)
  \end{tikzpicture}
\end{center}
\end{tkzexample}

\subsubsection{方法 \tkzMeth{point}{north(d)}}
\label{ssub:example_method_point_north_d}

该方法创建一个新点，它位于过给定点、竖直向上（朝向北方）的直线上，与给定点相距一段竖直距离。

\medskip

当你想在某个参考点上方偏移特定距离处构造一个点时——例如放置一个标签，或构造一个已知高度的几何构型——它特别有用。

\medskip

可选参数 \code{d} 表示竖直距离。若省略，则使用默认值 $1$。

\vspace{1em}


\begin{tkzexample}[latex = 7cm]
\directlua{
 init_elements()
 z.O = point(0, 0)
 z.A = z.O:east()
 z.Ap= z.O:east(1):north(1)
 z.B = z.O:north()
 z.C = z.O:west()
 z.D = z.O:south()}
\begin{center}
\begin{tikzpicture}
   \tkzGetNodes
   \tkzDrawPolygon(A,B,C,D)
   \tkzDrawPoints(A,B,C,D,O,A')
\end{tikzpicture}
\end{center}
\end{tkzexample}

\subsubsection{方法 \tkzMeth{point}{normalize()}}
\label{ssub:method_normalize}

该方法返回一个新点，它位于从原点到当前点的线段上，与原点相距 $1$。它通常用于提取一个向量的方向，并将其长度归一化为 1。

\medskip

你也可以用该方法，沿给定方向构造一个与另一点相距固定距离的点。例如：

\begin{center}
\code{z.U = (z.C - z.B):normalize() + z.B}
\end{center}


这里，向量 $\overrightarrow{BU}$ 的长度为 $1$，且 $U$ 位于线段 $[BC]$ 上、从 $B$ 指向 $C$ 的方向。

\medskip

有两种等价方式可以得到相同结果：

\begin{mybox}
\begin{verbatim}
  z.U = z.C:normalize_from(z.B)
  z.U = L.BC:normalize()
\end{verbatim}
\end{mybox}


第二种方式需要事先创建直线对象 \code{L.BC}。

\vspace{1em}

\begin{tkzexample}[latex=8cm]
\directlua{
  init_elements()
  z.A  = point(0, 0)
  z.B  = point(4, 3)
  z.C  = point(1, 5)
  L.AB = line(z.A, z.B)
  L.BC = line(z.B, z.C)
  z.N = z.B:normalize()
  z.U = z.C:normalize_from(z.B)}
\begin{tikzpicture}
  \tkzGetNodes
  \tkzDrawLines(A,B  B,C)
  \tkzDrawPoints(A,B,C,U,N)
  \tkzLabelPoints(A,B,C,U,N)
  \tkzDrawSegment(B,C)
\end{tikzpicture}
\end{tkzexample}

\subsubsection{方法 \tkzMeth{point}{orthogonal(d)}}
\label{ssub:orthogonal_method}

设 $O$ 为平面原点，$A$ 为异于 $O$ 的一点。该方法构造一个新点 $B$，使得向量 $\overrightarrow{OB}$ 与 $\overrightarrow{OA}$ 正交：

\[
\overrightarrow{OB} \perp \overrightarrow{OA}
\]

\medskip
\DeleteShortVerb{\|}
默认地，选取点 $B$ 使 $|OB| = |OA|$。若给出可选参数 \code{d}，则构造点 $B$ 使 $|OB| = d$。
\MakeShortVerb{\|}

\medskip

该方法有助于构造垂直向量，或在与给定方向正交的圆上生成点。

\vspace{1em}


\begin{tkzexample}[latex=7cm]
\directlua{
  init_elements()
  z.A = point(3, 1)
  z.B = z.A:orthogonal(1)
  z.O = point(0, 0)
  z.C = z.A:orthogonal()}
\begin{center}
  \begin{tikzpicture}[gridded]
    \tkzGetNodes
    \tkzDrawSegments(O,A O,C)
    \tkzDrawPoints(O,A,B,C)
    \tkzLabelPoints[below right](O,A,B,C)
  \end{tikzpicture}
\end{center}
\end{tkzexample}

\subsubsection{方法 \tkzMeth{point}{at(pt)}}
\label{ssub:_point_at_method}

该方法是对 \code{orthogonal} 方法的补充。它不是构造使 $\overrightarrow{OB} \perp \overrightarrow{OA}$（以 $O$ 为原点）的点 $B$，而是构造一个满足下式的点 $B$：

\[
\overrightarrow{AB} \perp \overrightarrow{OA}
\]


在此情形下，参考方向仍为 $\overrightarrow{OA}$，但正交向量是从点 $A$ 出发构造的，而非从原点出发。结果得到点 $B$，它位于过 $A$ 且与直线 $(OA)$ 正交的直线上。

\medskip

当处理局部正交方向时，该方法很有用，例如在三角形中作高，或定义锚定于给定点的垂直向量。

\vspace{1em}

\begin{tkzexample}[latex=7cm]
 \directlua{
 init_elements()
 z.O = point(0, 0)
 z.A = point(3, 2)
 z.B = z.A:orthogonal(1)
 z.C = z.A + z.B
 z.D =(z.C - z.A):orthogonal(2):at(z.C)}
 \begin{center}
   \begin{tikzpicture}[gridded]
   \tkzGetNodes
   \tkzLabelPoints[below right](O,A,C)
   \tkzLabelPoints[above](B,D)
   \tkzDrawSegments(O,A A,B A,C C,D O,B)
   \tkzDrawPoints(O,A,B,C,D)
   \end{tikzpicture}
 \end{center}
\end{tkzexample}

%====================================================================
\subsubsection{方法 \tkzMeth{point}{PPP(a,b)}}
\label{ssub:point_PPP}
%====================================================================
该方法在文档 \tkzname{Geometry Euclidean} [\href{https://altermundus.fr}{AlterMundus}] 中介绍。

\begin{tkzexample}[latex=.4\textwidth]
\directlua{
  init_elements()
  z.A = point(0, 0)
  z.B = point(4, 0)
  z.C = point(3, 4)
  PA.center, PA.through, _ = z.A:PPP(z.B, z.C)
  tkz.nodes_from_paths(PA.center, PA.through)
}

\begin{center}
  \begin{tikzpicture}
    \tkzGetNodes
    \tkzDrawPolygon(A,B,C)
    \tkzDrawCircles(w1,t1)
    \tkzDrawPoints(A,B,C)
  \end{tikzpicture}
\end{center}
\end{tkzexample}



\subsubsection{方法 \tkzMeth{point}{shift\_orthogonal\_to(pt, dist)}}
\label{shift_orthogonal}

\medskip
\texttt{语法： } \verb|z.P = z.A:shift_orthogonal_to(z.B, 2)|

\medskip
\texttt{返回值}
该方法返回将当前点 A 沿与直线 (AB) 正交的方向平移一个有向距离 dist 后得到的点。
正距离对应于将向量 $\overrightarrow{AB}$ 绕点 A 旋转 $+90^\circ$，
负距离对应于旋转 $-90^\circ$。

\begin{verbatim}
z.Q = z.B:shift_orthogonal_to(z.A, 2)
\end{verbatim}
正距离对应于将向量 $\overrightarrow{BA}$ 绕点 B 旋转 $+90^\circ$。

\begin{tkzexample}[latex=.4\textwidth]
\directlua{
z.A = point(3, 2)
z.a = point(3, 3)
z.B = point(7, 3)
z.b = point(7, 4)
C.A = circle(z.A, z.a)
C.B = circle(z.B, z.b)
z.TA = z.A:shift_orthogonal_to(z.B, C.A.radius)
z.TB = z.B:shift_orthogonal_to(z.A, -C.A.radius)
}
\begin{tikzpicture}[gridded]
\tkzGetNodes
\tkzDrawCircles(A,a B,b)
\tkzDrawLines(A,B TA,TB)
\tkzDrawPoints(A,B,TA,TB)
\end{tikzpicture}
\end{tkzexample}

\subsubsection{方法 \tkzMeth{point}{shift\_collinear\_to(pt, dist)}}
\label{shift_collinear}
该方法返回将当前点 A 沿与直线 (AB) 相同的方向平移一个有向距离 dist 后得到的点。

\medskip
\texttt{语法： } \verb|z.P = z.A:shift_collinear_to(z.B, 2)|



\subsubsection{方法 \tkzMeth{point}{rotation(an, obj)}}
\label{ssub:object_rotation}

与后面几个方法一样，绕一点的旋转可应用于一个点、一组点或一个对象。

该方法绕当前点（作为旋转中心）旋转一个或多个点或一个对象。


\medskip
\texttt{用途：}

该方法将一个几何对象绕指定点旋转给定角度。

在本例中，三角形绕作为旋转中心的点 $O$ 旋转角度 $\pi/3$。


\medskip
\texttt{语法： } \verb|new_obj = z.I:rotation(an, obj)|


\medskip
\texttt{返回：}
\begin{itemize}
\item 结果是一个点（或一组点），由每个给定点绕中心旋转指定角度得到。
\item 该方法返回一个同类型的新对象，并作相应旋转。
\end{itemize}


\medskip
\texttt{对一个点或一组点的用法示例}

在下面的示例中，一组点绕给定中心旋转。

\vspace{1em}

\begin{tkzexample}[latex=7cm]
\directlua{
 init_elements()
 z.a = point(0, -1)
 z.b = point(4, 0)
 z.o = point(6, -2)
 z.ap,
 z.bp = z.o:rotation(math.pi / 2, z.a, z.b)}
\begin{center}
\begin{tikzpicture}[scale = .5]
 \tkzGetNodes
 \tkzDrawLines(o,a o,a' o,b o,b')
 \tkzDrawPoints(a,a',b,b',o)
 \tkzLabelPoints(b,b',o)
 \tkzLabelPoints[below left](a,a')
 \tkzDrawArc(o,a)(a')
 \tkzDrawArc(o,b)(b')
\end{tikzpicture}
\end{center}
\end{tkzexample}

\medskip
\texttt{对一个对象的用法示例}

\vspace{1em}
\begin{tkzexample}[latex=7cm]
\directlua{
 init_elements()
 z.O  = point(-1, -1)
 z.A  = point(2, 0)
 z.B  = point(5, 0)
 L.AB  = line(z.A, z.B)
 T.ABC = L.AB:equilateral()
 S.fig = L.AB:square()
 _,_,z.E,z.F = S.fig:get()
 S.new = z.O:rotation(math.pi / 3, S.fig)
 _,_,z.Ep,z.Fp = S.new:get()
 z.C = T.ABC.pc
 T.ApBpCp = z.O:rotation(math.pi / 3, T.ABC)
 z.Ap,z.Bp,
 z.Cp = T.ApBpCp:get()}
\begin{center}
\begin{tikzpicture}[scale = .6]
 \tkzGetNodes
 \tkzDrawPolygons(A,B,C A',B',C' A,B,E,F)
 \tkzDrawPolygons(A',B',E',F')
 \tkzDrawPoints(A,B,C,A',B',C',O)
 \tkzLabelPoints(A,B,C,A',B',C',O)
 \begin{scope}
  \tkzDrawArc[delta=0,->,dashed,red](O,A)(A')
  \tkzDrawSegments[dashed,red](O,A O,A')
 \end{scope}
\end{tikzpicture}
\end{center}
\end{tkzexample}




\subsubsection{方法 \tkzMeth{point}{symmetry(obj)}}
\label{ssub:object_symmetry}

\medskip
\texttt{描述：}
该方法关于给定点对一个几何对象作中心对称（点反射）。

\medskip
\texttt{参数}

参数 \code{obj} 可以是点、直线、三角形、圆或任何其他受支持的几何对象。该对象的每个元素都关于中心点作反射，从而产生一个同类型的新对象。

\medskip
\texttt{用法示例}

下面的示例展示如何以一个参考点为中心，对一个对象（例如三角形）施加中心对称。

\vspace{1em}

\begin{tkzexample}[latex=7cm]
\directlua{
  init_elements()
  z.a = point(0, -1)
  z.b = point(2, 0)
  L.ab = line(z.a, z.b)
  C.ab = circle(z.a, z.b)
  z.o = point(1, 1)
  z.ap, z.bp = z.o:symmetry(C.ab):get()}
\begin{center}
  \begin{tikzpicture}
  \tkzGetNodes
  \tkzDrawCircles(a,b a',b')
  \tkzDrawLines(a,a' b,b')
  \tkzDrawLines[red](a,b a',b')
  \tkzDrawPoints(a,a',b,b',o)
  \tkzLabelPoints(a,a',b,b',o)
  \end{tikzpicture}
\end{center}
\end{tkzexample}

\subsubsection{方法 \tkzMeth{point}{homothety(k, obj)}}
\label{ssub:method_point_homothety_k_obj}

\texttt{用途： } 该方法关于当前点（作为变换中心）对一个几何对象作位似（放大或缩小）。

\medskip

\texttt{参数： }
\begin{itemize}
  \item \code{k} — 位似比（一个实数），
  \item \code{obj} — 待变换的对象，可以是：
  \begin{enumerate}
    \item 单个点，
    \item 一组点，
    \item 或一个几何对象（直线、三角形、圆等）。
  \end{enumerate}
\end{itemize}


正的比 \code{k} 产生以该点为中心的缩放，而负的比还会使对象关于中心作反射。

\vspace{1em}

\begin{tkzexample}[latex=7cm]
\directlua{
 init_elements()
 z.A = point(0, 0)
 z.B = point(1, 2)
 z.E = point(-3, 2)
 z.C, z.D = z.E:homothety(2, z.A, z.B)}
\begin{center}
\begin{tikzpicture}[scale = .5]
  \tkzGetNodes
  \tkzDrawPoints(A,B,C,E,D)
  \tkzLabelPoints(A,B,C,E)
  \tkzDrawCircles(A,B C,D)
  \tkzDrawLines(E,C E,D)
\end{tikzpicture}
\end{center}
\end{tkzexample}

 该方法将一个点的坐标转换为格式化字符串，可直接在正文中显示。

\medskip

小数位数由全局变量 \code{tkz\_dc} 控制，在 \code{init\_elements()} 函数中默认设为 $2$。你可以在调用 \code{print()} 之前赋予新值来覆盖它：

\begin{verbatim}
tkz_dc = 0
\end{verbatim}

\medskip

当在数学表达式中显示坐标、和、积或中间结果时，这特别有用。

\medskip
\texttt{用法示例：}

\begin{verbatim}
\directlua{
  init_elements()
  z.A = point(1, 2)
  z.B = point(1, -1)
  z.a = z.A + z.B
  z.m = z.A * z.B
  tkz_dc = 0}
  点 $A$ 和 $B$ 对应的复数分别为
   \tkzUseLua{z.A:print()} 与 \tkzUseLua{z.B:print()}，
  它们的和为 \tkzUseLua{z.a:print()}，
  它们的乘积为 \tkzUseLua{z.m:print()}。
\end{verbatim}

\directlua{
init_elements()
z.A = point(1, 2)
z.B = point(1, -1)
z.a = z.A + z.B
z.m = z.A * z.B
tkz_dc = 0
}

点 $A$ 和 $B$ 的各自附标为
   \tkzUseLua{z.A:print()} 和 \tkzUseLua{z.B:print()}，
 它们的和为 \tkzUseLua{z.a:print()}，
 积为 \tkzUseLua{z.m:print()}。

  %%%%%%%%%%%%%%%%%%%%%%%%%%%%%%%%
  %%%  实用工具 %%%
  %%%%%%%%%%%%%%%%%%%%%%%%%%%%%%%%
  \subsubsection{\tkzMeth{point}{identity}(pt)}
  \label{ssub:identity}

  \medskip
  \texttt{语法： } \verb|z.A:identity(z.B)|

  \medskip
  \texttt{用途： }判断两个点在 \tkzVar{tkz}{tkz.epsilon} 所定义的数值精度内是否相同。

  \medskip
  \texttt{参数： } 要与 \textbf{self} 比较的点。\\

  \medskip
  \texttt{返回值： } 若两个点在数值上相同则为 True，否则为 false。\\


  \begin{tkzexample}[latex=.5\textwidth]
  \directlua{
   z.A = point(0, 0)
   z.B = point(0.0000000001, 0)
   tex.print(tostring(z.A:identity(z.B)))
   tex.print('\\\\')
   tex.print(tkz.epsilon)}
   %% true if tkz.epsilon > 1e-10
  \end{tkzexample}

  \medskip
  \texttt{说明： }

  该方法通过计算两个点之间的欧氏距离并检查其是否小于全局容差 \textbf{tkz.epsilon} 来比较它们。
  在浮点计算中检验点是否相等时它很有用。



  \subsubsection{\tkzMeth{point}{print}()}
  \label{ssub:print}

  \texttt{用法：}
  \begin{verbatim}
  z.A:print()
  \end{verbatim}

  \medskip

  该方法以适合 \LaTeX{} 的格式打印该点的坐标。

  它在内部使用 \texttt{tostring(self)} 将该点转换为字符串，并通过 \texttt{tex.print} 将结果发送到 \TeX{} 输出。因此输出可直接用于文档。

  \medskip

  \texttt{返回：} 无。

  \medskip

  \texttt{示例：}
  \begin{verbatim}
  \directlua{
    z.A = point(1, 2)
  }
  Coordinates of A: \tkzUseLua{z.A:print()}
  \end{verbatim}

  \medskip

  \texttt{说明：}

  输出的格式取决于 \texttt{point} 类的 \texttt{\_\_tostring} 方法。通常它会产生形如 \texttt{(x,y)} 的坐标，可直接用于 TikZ 路径。
\endinput
